\section{PG-GAT v2 --- Visual Feature Augmentation}
\label{sec_meth_v2}

This section describes the visual-feature-augmented variant of PG-GAT.
The variant tests the hypothesis that pre-trained image features sampled at each detected keypoint
add identity signal that the positional input of Section~\ref{sec_meth_problem_graph} cannot provide.
The methodology is recorded here so that the Chapter~\ref{chap_4} negative result is interpretable;
the discussion of why the visual signal does not translate into a useful headline gain belongs to Section~\ref{sec_disc_v2}.

\subsection{Motivation and design constraints}
\label{sec_meth_v2_motivation}

The PG-GAT embedding of Section~\ref{sec_meth_pggat} consumes only normalised keypoint positions and joint types.
The image itself is consumed only by the upstream keypoint detector,
and no per-pixel visual context is propagated to the grouping stage.
A complementary lightweight feasibility study (\texttt{eval\_mobilenet\_features.py})
sampled MobileNetV2~\cite{sandler2018mobilenetv2} features at each detected keypoint
and showed that the addition of those features to the input of an otherwise unchanged $k$-means improves COCO grouping by a positive but diminishing amount as the grouping method becomes structurally stronger:
six points on plain $k$-means,
three points on type-constrained $k$-means,
and less than one point on Hungarian-matched type-wise assignment.
That diminishing trajectory is a methodological warning sign:
the strongest grouping methods barely use the visual signal,
which already places a ceiling on what a trained variant can extract from it.
The methodology below was designed to give a trained PG-GAT v2 every advantage the feasibility study did not,
so that the resulting negative result reported in Chapter~\ref{chap_4} cannot be attributed to an underpowered architecture.

\subsection{Visual feature extraction and caching}
\label{sec_meth_v2_features}

The feature backbone is MobileNetV2~\cite{sandler2018mobilenetv2} pre-trained on ImageNet~\cite{deng2009imagenet},
specifically the output of layer~17,
which produces a $320$-channel feature map at $1/32$ image resolution.
The choice of layer is methodological:
shallower layers (mid-level textures and edges) were shown by the feasibility study to hurt every grouping method tested,
and only the deepest semantic features produced any gain.
The backbone is held frozen throughout;
no end-to-end joint training of the backbone with PG-GAT v2 is attempted,
which keeps the comparison against PG-GAT controlled for everything except the visual input.

For each COCO image,
the backbone is run once and the resulting feature map is bilinearly sampled at every ground-truth keypoint pixel location;
the result is a per-keypoint feature vector $\mathbf{u}_i \in \mathbb{R}^{320}$.
This sampling is performed once over the COCO \texttt{train2017} and \texttt{val2017} splits and the per-image cache is written to disk
as a PyTorch Geometric \texttt{Data} object with the feature attribute attached at the node level.
The cached dataset contains $54{,}564$ training samples and $2{,}258$ validation samples,
and removes the backbone forward pass from the training loop entirely.
Training on cached features takes a fraction of the time that on-the-fly extraction would,
and trades memory for the ability to iterate quickly over PG-GAT v2 design choices.

\subsection{Architecture and input fusion}
\label{sec_meth_v2_arch}

The PG-GAT v2 network extends the baseline PG-GAT architecture of Section~\ref{sec_meth_pggat}
with a visual-feature projection block at the input stage.
Each cached MobileNet feature vector $\mathbf{u}_i$ is projected through
a learned linear layer to dimension $32$,
followed by a layer normalisation~\cite{ba2016layernorm} and a GELU activation~\cite{hendrycks2016gelu}:
\begin{equation}
    \mathbf{u}_i^{\text{proj}}
    = \mathrm{GELU}\!\bigl(\mathrm{LayerNorm}(\mathbf{W}_v \mathbf{u}_i + \mathbf{b}_v)\bigr),
    \quad \mathbf{W}_v \in \mathbb{R}^{32 \times 320}.
    \label{eq:v2_projection}
\end{equation}
The projected vector $\mathbf{u}_i^{\text{proj}} \in \mathbb{R}^{32}$ is concatenated with the positional and joint-type features of Equation~\eqref{eq:node_feature_no_depth}
before the first GAT layer.
All three skeleton-aware modifications of Section~\ref{sec_meth_pggat} are retained unchanged;
visual augmentation is therefore strictly additive to the PG-GAT design rather than a replacement for any of its components.

The total parameter count of PG-GAT v2 is $197{,}280$,
of which approximately $27{,}000$ are added by the projection block of Equation~\eqref{eq:v2_projection}
and the small adjustment to the first GAT layer's input dimension.
The remaining $170{,}560$ parameters are inherited from PG-GAT.
The MobileNet backbone itself contributes a further $2.2 \times 10^{6}$ parameters
which are not optimised but do count toward any apples-to-apples comparison against HigherHRNet's $28.6 \times 10^{6}$;
Section~\ref{sec_results_model_size} addresses the model-size accounting in full.

\subsection{Training schedule and the extended run}
\label{sec_meth_v2_training}

PG-GAT v2 is trained from scratch on the cached COCO training features
with the standard contrastive loss of Equation~\eqref{eq:contrastive_loss}
and no grouping head.
Two training schedules are reported in Chapter~\ref{chap_4}:
a 20-epoch run at learning rate $10^{-3}$ with cosine decay to $10^{-5}$,
and a 50-epoch extended run that resumes from the 20-epoch checkpoint at a halved initial learning rate $5 \times 10^{-4}$
and continues for 30 additional epochs.
The extended schedule is included to control against the possibility that a short training run undersells the architecture's capacity to extract value from the visual features;
the result reported in Chapter~\ref{chap_4} shows that the additional 30 epochs add at most a few thousandths of a PGA point above the 20-epoch number,
which establishes that the comparison against PG-GAT is not limited by undertraining.

The optimiser and clipping settings are inherited from Section~\ref{sec_meth_baseline_training}.
Validation runs every two epochs on the cached COCO validation features,
and the checkpoint with the best COCO validation PGA is retained as the final model.
The decision to validate on COCO rather than on synthetic data is methodological:
the synthetic generator does not produce per-keypoint MobileNet features,
and constructing them from rendered images would require running the backbone on a domain it was not trained for
and would introduce a confounding domain gap.
Reporting only COCO validation isolates the architectural question without that confound;
a footnote in Section~\ref{sec_results_v2} reports the synthetic transfer separately for completeness.
